An accurate trigonometric parallax calibration for Galactic Cepheids has long
been sought, but is very difficult to achieve in practice. All known classical
(Galactic) Cepheids are more than 250\,pc away, therefore for direct distance
estimates to achieve an uncertainty of up to 10\%, parallax uncertainties of
up to $\pm 0.2$\,milliarcsec are needed, requiring space-based observations.
The Hipparcos satellite reported parallaxes for 200 of the nearest Cepheids,
but even the best of these had high uncertainties. Recent progress has come
with the use of the Fine Guidance Sensor on HST where parallaxes (in many
cases) accurate to better than $\pm 10$\% were obtained for 10 Cepheids,
spanning a range of periods from 3.7 to 35.6 days. These nearby Cepheids cover
distances from about 300 to 560\,pc.

The measured periods, $P$, and average magnitudes in $V$, $K$ and $I$ bands
are given in Table~1 as well as the $A_V$ and $A_K$ for extinction in $V$ and
$K$ bands, respectively. The measured parallaxes with their uncertainties are
also given in milliarcsec (mas). All measured apparent magnitudes have
negligible uncertainty.

Table 1: Periods and average apparent magnitudes of 5 Galactic Cepheids with
accurate parallax measurements.

\begin{center}
\begin{tabular}{l|c|c|c|c|c|c|c|c}
 & $P$ & $\langle V\rangle$ & $\langle K\rangle$ & $A_V$ & $A_K$
 & $\langle I\rangle$ & parallax & error \\
 & (day) & (mag) & (mag) & (mag) & (mag) & (mag) & (mas) & (mas) \\
\hline
RT Aur      & 3.728  & 5.464 & 3.925 & 0.20 & 0.02 & 4.778 & 2.40 & 0.19 \\
FF Aql      & 4.471  & 5.372 & 3.465 & 0.64 & 0.08 & 4.510 & 2.81 & 0.18 \\
X Sgr       & 7.013  & 4.556 & 2.557 & 0.58 & 0.07 & 3.661 & 3.00 & 0.18 \\
$\zeta$ Gem & 10.151 & 3.911 & 2.097 & 0.06 & 0.01 & 3.085 & 2.78 & 0.18 \\
l Car       & 35.551 & 3.732 & 1.071 & 0.52 & 0.06 & 2.557 & 2.01 & 0.20 \\
\end{tabular}
\end{center}

\begin{description}
\item[(D1.1)] The observed correlation between the period of a Cepheid and its
  brightness is usually described by the so-called ``Period-Luminosity (PL)
  relation'', where $L \propto P^\beta$. In fact, such a relation is normally
  expressed in terms of the period and absolute magnitude, instead of
  luminosity. Hereafter, we shall refer to the Period-Absolute magnitude
  relation as the conventionally named ``PL relation''.

  Use the data given in Table~1 to plot a suitable linear graph in order to
  derive the Cepheid PL relation for the $V$-band and $K$-band. You should
  plot each graph separately on different pieces of graph paper. Determine the
  slope of the line that best describes the linear relation of the data. (You
  may find the relation
  $\Delta(\log_{10} x) \approx \dfrac{\Delta x}{x \log_e 10}$ useful)
  \hfill[36.5 Marks]
\end{description}

Any apparent differences in PL relations of stars in the different bands can
be explained if one also considers differences in colour. Therefore, the PL
relation is in fact a PLC (Period-Luminosity-Colour) relation. This is from
the reddening effect, due to extinction being a function of wavelength, which
can also vary among different Cepheids due to their different metallicities,
foreground Interstellar Medium and dust.

A new reddening-free magnitude (or bandpass) called ``Wesenheit'' has been
proposed that does not require the explicit information of the extinction of
individual stars but uses colour information from the star itself to get rid
of the effect. For example, $W_{VI}$ use $V$ and $I$ band photometry and is
defined as
\[ W_{VI} = V - \left[\frac{A_V}{E(V-I)}\right](V-I), \]
\[ \phantom{W_{VI}} = V - R_v (V-I) \]
where $R_v$ depends on the reddening law. In this case, we shall take the
value of $R_v$ to be 2.45.

\begin{description}
\item[(D1.2)] From the data given in Table~1, plot and derive the
  reddening-free PL relation using Wesenheit $W_{VI}$ magnitudes. Estimate the
  linear slope of the relation as well as its uncertainty.
  \hfill[14.5 Marks]
\item[(D1.3)] Next, we would like to use the newly-derived PL relations from
  question (D1.1) \& (D1.2) to estimate the distance to the Large Magellenic
  Cloud (LMC) using periods and magnitudes of classical Cepheids in the LMC.
  In Table~2, the periods, average extinction-corrected apparent magnitudes,
  $\langle V_{\mathrm{corr}}\rangle$, and Wesenheit $W_{VI}$ magnitudes are
  given.

  Estimate the distance modulus, $\mu$, to each star and then use all the
  information to derive the distance to the LMC (in parsecs) and its standard
  deviation for each band.

  Compare if the derived distances are statistically different for the 2 bands
  (YES/NO).

  Are the standard deviations of the estimated distances for 2 bands different
  (YES/NO)?

  Based on this dataset, which band ($V$ or Wesenheit) is more accurate in
  estimating the distance to the LMC? \hfill[24 Marks]
\end{description}

Table 2: Period, average extinction-corrected apparent magnitude,
$\langle V_{\mathrm{corr}}\rangle$, and average Wesenheit magnitude
measurements of Cepheids in the LMC.

\begin{center}
\begin{tabular}{l|c|c|c}
 & $P$ & $\langle V_{\mathrm{corr}}\rangle$ & $\langle W_{VI}\rangle$ \\
 & (day) & mag & mag \\
\hline
HV12199 & 2.63  & 16.08 & 14.56 \\
HV12203 & 2.95  & 15.93 & 14.40 \\
HV12816 & 9.10  & 14.30 & 12.80 \\
HV899   & 30.90 & 13.07 & 10.97 \\
HV2257  & 39.36 & 12.86 & 10.54 \\
\end{tabular}
\end{center}
